\subsection{Articulación entre prueba de entrega, DTE y acuse}
\label{sec:pod-dte-acuse}

La guía de despacho electrónica debe estar vinculada a la salida de mercadería antes de iniciar el traslado. La aplicación no la emite: M5 entrega los datos de despacho por la capa anticorrupción al ERP, que conserva la responsabilidad tributaria y devuelve folio, estado y documento. M6 captura después la recepción efectiva, incluidos bultos rechazados, identidad del receptor y evidencia. El acuse técnico de recepción o procesamiento ante el SII y el acuse de recepción de mercaderías del destinatario son hechos distintos; cada uno conserva su origen, fecha y estado.

\subsubsection{Secuencia y excepciones}

La secuencia lógica es: preparación confirmada en M5; solicitud de guía al ERP mediante ACL; autorización de salida con el documento disponible; captura de POD y diferencias en M6; conciliación de la entrega en M7/M8; y, si corresponde, solicitud al ERP de factura, nota de crédito o tratamiento tributario de la devolución. Un mismo \texttt{transaction\_id} relaciona pedido, guía, entrega, evidencia, acuses y documentos posteriores. El acuse del destinatario se obtiene por el mecanismo que acuerde el CLIENTE conforme a la normativa aplicable; una fotografía o un QR no se presume equivalente por sí solo.

Si el ERP o su canal tributario no confirma la guía antes de la salida, M5 muestra el despacho como bloqueado y lo deriva al responsable de Operaciones y Tributación para aplicar únicamente la contingencia autorizada por el CLIENTE. La captura offline del POD permite registrar una entrega ya despachada; no autoriza emitir retroactivamente una guía desde el teléfono. Esta separación protege la continuidad de la evidencia sin atribuir a la aplicación facultades tributarias que permanecen en el ERP.

\subsubsection{Control de liberación en la ventana crítica}

M5 conserva los estados carga confirmada, guía solicitada, resultado incierto, guía disponible y salida liberada. La solicitud lleva una clave estable por despacho y versión de carga. Ante un timeout se consulta el resultado de esa misma solicitud antes de repetirla; no se genera otra guía a ciegas. Solo se libera cuando la guía corresponde a la carga vigente y su evidencia está disponible localmente. Una modificación posterior de carga vuelve a exigir la validación documental del ERP. El acuse técnico del SII, el documento emitido y la autorización de traslado no se representan con un único booleano.

Al cerrar la carga nocturna se solicita la guía por adelantado, sin confundir anticipación con contingencia. Antes de las 05:30, Operaciones revisa las salidas programadas y las guías disponibles. Durante 05:30--07:00 se prioriza INT-07 sobre reportes y cargas masivas; una incidencia identifica los camiones afectados, hora prevista, responsable y estado documental. Los despachos con documentación válida siguen su curso; los afectados se retienen hasta recuperar el ERP o aplicar el procedimiento tributario aprobado. Ningún perfil de supervisor puede omitir ese control con una autorización genérica.

Esta retención protege la validez del despacho, pero puede incumplir la ventana de indisponibilidad cero del caso. La solución exige cerrar AL-DTE-01: procedimiento firmado por Operaciones y Tributación del CLIENTE, compatible con su emisor ERP, y ensayo de las 96 salidas con fallas inyectadas. No se atribuye aprobación a ese procedimiento mientras no exista el acta. La preemisión reduce exposición, pero no resuelve cambios de carga ni fallas prolongadas por sí sola.

El Anexo 4.1-U distingue emisión tributaria, constancia operativa y acuse del receptor. La firma, fotografía o QR del POD no se equiparan automáticamente al recibo exigible: el mecanismo se selecciona según el receptor y se canaliza por el ERP o el perfil EDI aprobado (SII, 2018, s.~f.-a, s.~f.-b).


